\section{Discussion}\label{sec:discussion}

\subsection{What the evidence so far supports}

It is worth being exact about which of the framework's predictions the paper has tested. None of P1--P6 is yet tested with the data it was written for. The cross-section of Section~\ref{sec:cross} shows that official values vary strongly with 2011 inundation and that most of the variation is removed by land cover measured ten years later, which may itself be an outcome of the flood; it neither confirms nor rejects P1, and it rules out the cross-sectional route to it. The design of Section~\ref{sec:design} is powered to test P2 in its first-round form and is waiting on the valuation calendar. The synthetic exercises of Sections~\ref{sec:credit} and~\ref{sec:stress} show that the estimators recover P3 and the breaking point when they hold by construction, which is a test of the code, not of the world. The $\lambda$-sensitivity of Section~\ref{sec:stress} is a property of the model that any calibration inherits: it says that P5's valuation term is quantitatively important, not that Thai valuations ignore hazard.

The descriptive evidence is consistent with the framework but does not discriminate between it and simpler stories. Credit in the most flooded provinces did not contract after 2011 but fell steadily behind the other provinces, appraisal prices fell relative to other regions for at least ten months, and the literature's own Thai study finds default below model-predicted levels after the flood \citep{chitchumnong2025physical}. Under the framework these are what a relief-censored default process and lagged valuations produce; they are also what a small true loss would produce. The project exists to tell the two apart.

\subsection{Threats and limits}

The threats to the pre-specified design are of four kinds, each with a registered response. \emph{Selection into flooding}: flooded sheets are lower, more agricultural and cheaper than dry ones; the design balances on these within province, caps control height, and reports the residual gaps on the unbalanced variables rather than suppressing them. \emph{Measurement error in treatment}: the radar maps have low recall against UNOSAT and imperfect precision against optical imagery; both push the estimate towards zero, the controls are required to be dry in every source, and the UNOSAT sample gives a second definition of treatment. \emph{Relief-driven revaluation}: committees might revise values in declared disaster districts for reasons other than flooding, which would make a negative estimate partly a relief effect; the declared-district sensitivity addresses it, at a large cost in controls. \emph{Changes in the schedule's structure between rounds}: re-drawn land units break the outcome; the pseudo-land-unit construction, the carry-over share and the Lee bounds handle attrition, and the stopping rules protect against a round that is a mechanical repricing of the previous one.

The deeper limit is the gap between $\lambda_T$ and $\lambda_B$. The official values are a tax base set by market comparison, and the central bank's collateral policy does not reference them. A negative estimate in Question~1 would show that the state's market-comparison valuer revises for a realised flood; it would not show that banks do. What it would supply is the first evidence from an emerging market on whether a market-comparison valuation of land responds to an observed flood, which is the ingredient the supervisory template takes for granted, and a benchmark against which bank appraisals, once observed, can be compared. The dependence of the breaking point on $\lambda_B$ in Table~\ref{tab:lambda} is the reason to obtain those appraisals.

Two further limits are honest statements of ignorance. The dispersion of value changes between rounds, $\sigma$, is unknown because no earlier round is public, and the power calculations bracket it rather than estimate it. And the 2011 satellite record cannot see the urban flood, so a 2011-based exposure measure for Bangkok understates the exposure of the largest part of the banking system; the World Bank's hazard-model mapping, not the satellite record, is the right input for Bangkok in the stress test.

\subsection{What the results would mean for policy}

The paper is not a policy document, but the quantities it is built to estimate map onto the choices in the public record. If official values do not respond to a realised flood in the first round after it, then a transfer-fee base and any valuation that leans on it carry stale hazard information for up to four years, and the 15\% within-round trigger is not doing that work. If flood exposure of collateral surfaces in stage~2 and restructuring rather than default, then the classification freeze in the standing relief template is doing what it is designed to do, and the expected-credit-loss framework, not the default statistics, is where supervisors would see the flood. And if the breaking point is as sensitive to the valuation parameter as the illustrative calibration suggests, then a physical-risk stress test that uses banks' own collateral values without asking whether they price hazard will overstate the distance to the threshold most for the loans that appear safest. These are conditional statements; the project's purpose is to remove the conditions.

\subsection{Next steps}

Five decisions precede registration and are recorded as open in the plan: the primary test for the radar samples, after the Monte Carlo found the placebo-footprint test rejecting a true null \mcxNullRi{} of the time while the wild cluster bootstrap held its size (Section~\ref{sec:montecarlo}; restricting the placebo pool did not repair the test and studentising it did, so the amended plan proposes the bootstrap as primary with a studentised placebo check); whether the 2021--23 radar reference period replaces the 2023-only reference as the primary treatment definition, which would drop \sensRefDroppedPct\% of treated sheets in exchange for a cleaner footprint; the overlap rule for Phayao, where \regControlsPhayao{} controls cannot be balanced on height and water; whether a pre-2011 land-cover product is obtained for the cross-section; and the public archive of the repository and sample files. Then, in order: register the pre-analysis plan before any valuation round is announced; obtain the 2024 and 2025 disaster declarations from the Department of Disaster Prevention and Mitigation to replace the stand-in lists in the declared-zone sensitivities; secure loan-level access through a formal, independent agreement, or fall back to listed-bank disclosures; run the registered analysis on the first published round, under whichever timing case materialises; and calibrate the reverse stress test with the estimates in hand. The analysis code for every step is written, tested on synthetic data and versioned.
